%!TEX root = ../../csuthesis_main.tex
\chapter{总结与展望}

本文主要任务是对车辆多体物理仿真系统进行优化，在CarlaUE4 仿真场景中针对仿真过程因动力学精度不足、多体约束误差、数据同步延迟及可视化耦合性差等问题导致的仿真失真，提出并实现了Chrono 与 CarlaUE4 联合仿真模型。该方案基于Chrono 多体动力学算法与CarlaUE4 虚幻引擎，通过引入高精度车辆子系统建模提升仿真平台对复杂工况、悬架运动及轮胎力学的求解能力，同时采用UDP 低延迟通信与预测迭代同步机制，以应对数据交互延迟、动力学与渲染不同步、实时性不足等挑战。上文各项实验数据表明，通过优化 Chrono 求解算法并深度耦合 CarlaUE4 可视化环境，从而实现的多体物理车辆仿真系统优化是有效且成功的。


\section{主要研究成果总结}

本研究以基于虚幻引擎（CarlaUE4）与 Chrono 算法的多体物理车辆仿真模型为核心对象，完成了模型构建、平台耦合、算法优化与实验验证全流程工作，有效解决了传统车辆多体物理仿真存在的场景还原度不足、动力学求解精度偏低、多体约束处理不精准、虚拟场景耦合性差等问题，构建出一套高精度、高实时性、高工程实用性的车辆联合仿真系统，主要研究成果如下：

\subsection{完成高精度车辆多体动力学模型搭建}

依托 Chrono::Vehicle 模块化、模板化建模体系，完成车架、悬架、转向、传动、轮胎等核心子系统的标准化建模，采用 JSON 参数化配置实现车辆模型快速实例化，确保多体约束关系与动力学求解精度达到工程级标准，可精准复现车辆真实运动特性。

\subsection{实现仿真模型关键技术优化}

针对多体约束精度不足、动力学求解效率与实时性矛盾、多平台数据交互延迟三大核心挑战，优化 Chrono 约束求解算法与自适应步长求解策略，提升动力学计算效率；统一双平台数据格式与仿真时钟，降低联合仿真同步误差，提升系统运行稳定性。

\subsection{通过对比实验验证模型性能提升}

以 IMU 三轴速度、加速度、角速度为核心评价指标，开展 Chrono 与 Carla 原生物理引擎（PhysX）对照实验。结果表明：常规行驶工况下，两引擎数据趋势基本一致；极端碰撞工况下，Chrono 在垂直动力学响应、悬架弹性阻尼特性、碰撞能量传递、车辆姿态变化等方面更贴合真实物理规律，仿真精度与真实性显著优于传统游戏级物理引擎。


\section{不足与未来期望}
本研究已完成多体物理车辆仿真模型的构建与优化，实现了预期研究目标，但受限于实验条件、模型复杂度与计算资源，模型仍存在进一步完善的空间，未来可从以下方向开展深入研究：


\subsection{提升模型参数精度与通用性}

基于实车测试数据对 Chrono 模型的轮胎侧偏特性、转向系统、碰撞恢复系数、悬架刚度阻尼等关键参数进行对标校准，缩小常规工况下与实车数据的偏差；建立标准化车辆参数库，拓展模型在不同车型、不同路面条件下的适配能力。同时在CarlaUE4联合chrono仿真时，车辆碰撞后会自动退出chrono回到phyx仿真。

\subsection{拓展复杂场景与多物理场耦合仿真}

引入雨雪、风沙、极端温度等环境因素，结合土壤力学、流体力学模块，实现车辆 — 环境全物理场耦合仿真；增加非铺装路面、越野地形、多车交互、人车混行等复杂场景，提升平台对自动驾驶长尾极端工况的覆盖能力。


\subsection{拓展工程应用与研究领域}

将模型应用于自动驾驶全栈算法测试、底盘性能标定、碰撞安全分析、智能驾驶训练等实际工程场景；进一步拓展至工程机械、轨道交通、无人机等多体系统仿真领域，构建通用化多体物理仿真平台。

未来将持续围绕高精度、高实时、高耦合的多体物理仿真技术开展研究，不断完善仿真模型性能，推动仿真平台从实验室研究走向实际工程应用，为智能交通与车辆智能化发展提供更坚实的技术支撑。
\newpage